\chapter{INTRODUCTION}

\section{Background}
Modern healthcare systems rely on accurate occupancy and flow data for clinical triage, resource allocation, and workflow management. In Ghanaian hospitals like Komfo Anokye Teaching Hospital (KATH) and KNUST Hospital, the Outpatient Department (OPD) receives most incoming patients \cite{osei2024}. Hundreds of individuals arrive daily, creating crowded waiting halls with mixed demographics: sick children, accompanying siblings, adult caregivers, visitors, and hospital staff. Nurses classify clinical acuity with the South African Triage Scale (SATS) and register patients in the national Lightwave Health Information Management System (LHIMS) \cite{addotey2023}.

Despite investments in digital record systems, initial waiting area headcount monitoring remains manual. Triage nurses and ward assistants scan crowded waiting areas on foot. They count waiting occupants by hand and verbally ask parents whether accompanying children are patients requiring care or healthy siblings. This manual counting process causes delays in prioritizing vulnerable arrivals. It skews emergency nurse-to-patient staffing ratios and introduces discrepancies into administrative records \cite{keles2023, javaid2024}.

Computer Vision (CV) and Deep Learning (DL) offer tools for non-invasive visual monitoring. Neural networks convert ceiling CCTV streams into spatial coordinate matrices. They detect persons, estimate demographic occupancy, and track corridor density in real time \cite{diop2025, jiang2022}. However, standard pedestrian detectors assume upright, unoccluded adults. Furthermore, computer vision models detect physical bodies. They cannot inherently determine whether a detected child is a patient or sibling, nor whether an adult is a caregiver or nurse. Standard off-the-shelf detectors fail to capture pediatric-specific body proportions, particularly in low-resource clinical settings.

\section{Motivation and Significance}
This study addresses the systematic undercounting and misclassification of pediatric persons in crowded healthcare waiting areas. Children under five years move unpredictably and stay close to their caregivers. Standard vision models treat all human bodies as uniform bounding boxes with a typical 1:3 width-to-height ratio. These generic models fail to capture morphological proportions that distinguish children from adults on embedded platforms \cite{lin2022}, or face features in non-frontal orientations \cite{liu2020}.

This work provides three key contributions:
\begin{enumerate}[(i)]
    \item \textbf{Mathematical and Biological Constraints:} Instead of relying solely on black-box models, this work applies discrete differential operators, rank-1 tensor separability, and cephalocaudal growth priors. These geometric constraints stabilize demographic classification in occluded scenes.
    \item \textbf{Operational Healthcare Telemetry:} Objective automated person counts provide near-real-time occupancy inference designed to interface with hospital information systems (such as LHIMS). Computer vision detects physical bodies rather than clinical illness. However, accurate demographic counts provide baseline crowding data to help administrators allocate triage staff and supplies during peak clinic hours.
    \item \textbf{Privacy-Preserving Edge Computing:} The framework runs on local low-power CPU hardware using AVX2 SIMD instructions. It outputs anonymous coordinate vectors without storing facial data or transmitting raw video to cloud servers. This design minimizes the exposure of directly identifying visual information \cite{milkaite2019}.
\end{enumerate}

\section{Problem Statement}
Standard detection algorithms often struggle to accurately count children in pediatric clinical triage due to three physical constraints:
\begin{enumerate}[(i)]
    \item \textbf{Small Patient Scale (Scale Problem):} Children occupy few pixels in the video frame. This low pixel count removes fine facial and body features, which can cause standard detectors to miss small bodies.
    \item \textbf{Poor Lighting (Contrast Problem):} Variable hospital lighting reduces visual contrast. Patients blend into background chairs and blankets, which can cause detectors to lose boundary edges.
    \item \textbf{Dynamic Movement (Tracking Problem):} Children move erratically. Fast patient movement can cause identity swaps and duplicate counts across video frames.
\end{enumerate}

These factors cause standard detectors to misclassify children as adults or miss them entirely. This creates a high Mean Absolute Percentage Error (MAPE). Consequently, general detectors provide unreliable demographic data for hospital triage.

\section{Aim and Objectives}

\subsection{General Objective}
The main goal is to implement and validate a privacy-preserving computer vision system for automated pediatric detection, allometric verification, and multi-object tracking.

\subsection{Specific Objectives}
To achieve this goal, we set six specific objectives:
\begin{enumerate}[(i)]
    \item Evaluate pixel-level edge computation using discrete image representations and separable gradient convolutions on clinical hardware.
    \item Evaluate how CCTV viewing angles and training data limitations affect demographic detection. Assess the improvement of a decoupled classification stage.
    \item Build a multi-stage pipeline that decouples generic person detection from targeted demographic classification.
    \item Formulate a cephalocaudal allometric metric ($R_{\text{ceph}} = L_{\text{torso}}/L_{\text{leg}}$) to separate observed children from adults using body proportions.
    \item Develop a weighted temporal filter to stabilize demographic classifications and prevent identity flips.
    \item Evaluate the system across clinical CCTV feeds to measure counting accuracy, tracking continuity, and edge CPU throughput.
\end{enumerate}

\section{Research Questions}
This work answers five research questions:
\begin{enumerate}[(i)]
    \item How does rank-1 separable convolution compare with 2D spatial convolution in arithmetic complexity on edge hardware?
    \item How do CCTV camera angles and data quality affect demographic detection? Does a decoupled classification stage improve precision?
    \item Does the skeletal cephalocaudal ratio ($R_{\text{ceph}} = L_{\text{torso}} / L_{\text{leg}}$) provide a reliable separation margin between children and adults under steep camera angles?
    \item How effectively does weighted temporal smoothing reduce posture noise and maintain tracklet identity?
    \item What are the trade-offs between native FP32 AVX2 SIMD and emulated FP16 execution for latency and throughput on standard CPUs?
\end{enumerate}

\section{Delimitation and Scope}
This study combines applied mathematics, computer vision, and healthcare operations. We evaluate data from ambulatory triage areas, waiting corridors, and consultation halls at KATH and the KNUST Hospital in Kumasi, Ghana.

The scope is limited to non-invasive optical detection, demographic classification, and tracking of pediatric and adult persons. In this study, the pediatric cohort operationally includes children under 12 years of age. The primary goal is counting children accurately despite severe CCTV perspective distortion. The computer vision model separates individuals by physical allometry. Pre-pubescent children exhibit large cephalocaudal ratios ($R_{\text{ceph}} \ge 0.80$) and shorter normalized stature ($\theta_{\text{adult}} = 0.32$). Post-pubescent adolescents who have attained adult height and limb proportions are grouped with the adult class by visual sensors. Ground-truth demographic labels were established by two independent observers who reviewed video recordings alongside clinic reception registers.

The system does not assess clinical acuity or determine patient status independently of hospital workflow records. True patient identification requires linking optical counts with registration or queue-zone management rules. We exclude facial recognition and biometric identity logging to protect patient privacy. All tracking and occupancy metrics rely solely on anonymous bounding box centroids and skeletal coordinates.

\section{Organization of the Thesis}
This thesis contains five chapters:
\begin{itemize}
    \item \textbf{Chapter 2 (Literature Review):} Traces object detection from hand-crafted features to deep neural networks. It reviews pediatric anthropometry and crowd obstruction handling. It defines the research gap.
    \item \textbf{Chapter 3 (Methodology):} Formulates the multi-stage pipeline, digital quantization, separable convolutions, CIoU regression, margin-expanded crop classification, allometric growth ratios, and temporal tracking.
    \item \textbf{Chapter 4 (Results):} Reports empirical benchmarks in clinical waiting areas and pharmacy counters, presents detection and counting metrics, and evaluates CPU acceleration benchmarks.
    \item \textbf{Chapter 5 (Conclusion):} Summarizes the main findings, describes practical limitations, provides deployment recommendations for hospital administrators, and outlines future research directions.
\end{itemize}

Chapter 2 reviews the existing literature on object detection, pedestrian tracking, and pediatric anthropometry to establish the theoretical background for this study.